\documentclass[11pt]{article}

% --------------------------------------------------------------------------
% Two ways to solve electromagnetic transients:
%   global implicit DAE (Dynawo)  vs.  nodal companion models (classical EMTP)
% Compile with: pdflatex EMT_methods.tex
% --------------------------------------------------------------------------

\usepackage[a4paper,margin=2.4cm]{geometry}
\usepackage{amsmath,amssymb}
\usepackage{booktabs}
\usepackage{array}
\usepackage{xcolor}
\usepackage[colorlinks=true,linkcolor=blue,urlcolor=blue]{hyperref}

\newcommand{\dd}{\mathrm{d}}
\newcommand{\hstep}{h}
\newcommand{\vv}{\mathbf{v}}
\newcommand{\ii}{\mathbf{i}}

\title{\bfseries Two ways to solve electromagnetic transients\\[2pt]
\large Global implicit DAE (Dynawo) versus nodal companion models (classical EMTP)}
\author{EMT exploration notes --- Dynawo}
\date{\today}

\begin{document}
\maketitle

\begin{abstract}
Electromagnetic-transient (EMT) simulation computes the \emph{instantaneous}
three-phase voltages and currents of a network of resistors, inductors,
capacitors, sources and switches. Two families of methods do this and they sit
at opposite ends of a design axis. The \textbf{classical EMTP} method
discretises \emph{each element locally} into a resistor plus a history current
source (a \emph{companion model}) and solves one real, linear nodal system per
step. The \textbf{global-DAE} method (as used by Dynawo) keeps each element as
its continuous physical equation and lets a general implicit integrator
discretise the \emph{whole symbolic system} at once, solving a nonlinear
residual by Newton iterations each step. This note derives both from the same
trapezoidal rule, makes the precise correspondence explicit, and tabulates the
trade-offs.
\end{abstract}

\tableofcontents

% ==========================================================================
\section{The problem}
% ==========================================================================

We want the instantaneous waveforms $v(t)$, $i(t)$ in every phase. The building
blocks obey, per phase,
\begin{equation}
\text{resistor: } v = R\,i,
\qquad
\text{inductor: } v = L\,\frac{\dd i}{\dd t},
\qquad
\text{capacitor: } i = C\,\frac{\dd v}{\dd t}.
\label{eq:constitutive}
\end{equation}
Together with Kirchhoff's laws (KCL/KVL) over the topology, these form a system
of differential and algebraic equations. The question is \emph{how to advance it
in time}. Both methods below use the same one-step integration rule, the
\textbf{trapezoidal rule},
\begin{equation}
y_{n+1} = y_n + \frac{\hstep}{2}\big(\dot y_n + \dot y_{n+1}\big),
\qquad \hstep = t_{n+1}-t_n,
\label{eq:trap}
\end{equation}
which is second-order accurate and A-stable. What differs is the
\emph{scope} at which \eqref{eq:trap} is applied.

% ==========================================================================
\section{Method A --- Global implicit DAE (Dynawo)}
% ==========================================================================

Collect all unknowns (node voltages, inductor currents, source currents, \dots)
into one vector $y$. The network is written as a single, generally implicit,
differential-algebraic system
\begin{equation}
F\big(y,\dot y,t\big) = 0,
\label{eq:dae}
\end{equation}
mixing differential equations (for the $L$ and $C$ states) and algebraic ones
(KCL at nodes, resistor laws, ideal-source laws). \textbf{No element is
pre-discretised}: an inductor contributes literally $v - L\,\dot i = 0$.

\paragraph{Discretisation lives in the solver.}
Solving \eqref{eq:trap} for the derivative gives
\begin{equation}
\dot y_{n+1} = \frac{2}{\hstep}\big(y_{n+1}-y_n\big) - \dot y_n
\;\equiv\; c_j\,y_{n+1} + \alpha_n,
\qquad
\boxed{\,c_j = \dfrac{2}{\hstep}\,},\quad
\alpha_n = -\frac{2}{\hstep}y_n-\dot y_n .
\label{eq:yp}
\end{equation}
Substituting \eqref{eq:yp} into \eqref{eq:dae} turns each step into a
\emph{nonlinear algebraic} system for the single unknown $y_{n+1}$,
\begin{equation}
\Phi(y_{n+1}) \;\equiv\;
F\!\left(y_{n+1},\; c_j\,y_{n+1}+\alpha_n,\; t_{n+1}\right) = 0,
\label{eq:residual}
\end{equation}
solved by Newton's method, $J\,\Delta y = -\Phi(y_{n+1})$, with the iteration
matrix
\begin{equation}
J = \frac{\partial F}{\partial y} + c_j\,\frac{\partial F}{\partial \dot y}.
\label{eq:jac}
\end{equation}
The whole network --- including every nonlinearity and switch --- is inside
\eqref{eq:residual}; the matrix $J$ is rebuilt/refactorised as needed but the
system is \emph{never} reduced to a constant admittance.

\paragraph{This is exactly Dynawo's \texttt{SolverTRAP}.}
In \texttt{DYNSolverTRAP.cpp}, \texttt{computeYP} implements \eqref{eq:yp}
verbatim,
\[
\texttt{vectorYp[i] = (2./h)*(yy[i] - vectorYSave[i]) - vectorYpSave[i];}
\]
and the nonlinear solve \eqref{eq:residual}--\eqref{eq:jac} is handed to KINSOL.
A worked single branch (source--$R$--$L$ loop, state $i$):
\begin{equation}
L\,\dot i = v_s(t) - R\,i
\;\;\xrightarrow{\eqref{eq:yp}}\;\;
L\Big(\tfrac{2}{\hstep}(i_{n+1}-i_n)-\dot i_n\Big) = v_s(t_{n+1}) - R\,i_{n+1},
\label{eq:Aexample}
\end{equation}
one (here linear) equation solved for $i_{n+1}$ each step.

% ==========================================================================
\section{Method B --- Nodal companion models (classical EMTP)}
% ==========================================================================

Here \eqref{eq:trap} is applied \emph{inside each storage element} before any
system assembly, replacing it by a \textbf{companion model}: a conductance in
parallel with a history current source that depends only on past values.

\paragraph{Inductor.} With $v=L\,\dd i/\dd t$ and the trapezoidal rule,
\begin{equation}
i_{n+1} = i_n + \frac{\hstep}{2L}\big(v_{n+1}+v_n\big)
= \underbrace{\frac{\hstep}{2L}}_{G_L}\,v_{n+1}
+ \underbrace{\Big(i_n + \frac{\hstep}{2L}v_n\Big)}_{I_{\mathrm{h},n}^{L}} ,
\qquad
\boxed{\,G_L=\dfrac{\hstep}{2L}\,}.
\label{eq:compL}
\end{equation}

\paragraph{Capacitor.} With $i=C\,\dd v/\dd t$,
\begin{equation}
i_{n+1} = \underbrace{\frac{2C}{\hstep}}_{G_C}\,v_{n+1}
\;-\;\underbrace{\Big(\frac{2C}{\hstep}v_n + i_n\Big)}_{-\,I_{\mathrm{h},n}^{C}},
\qquad
\boxed{\,G_C=\dfrac{2C}{\hstep}\,}.
\label{eq:compC}
\end{equation}

\paragraph{Resistor.} Simply $G_R = 1/R$, no history term.

\medskip
Every element is now \emph{resistor $\parallel$ known current source}. Each phase
node obeys KCL, and assembling them gives one real, symmetric, \textbf{constant}
(for fixed $\hstep$ and topology) conductance system
\begin{equation}
\boxed{\;\mathbf{G}\,\vv_{n+1} = \ii^{\,\text{src}}_{n+1} + I_{\mathrm{h},n}\;}
\label{eq:nodal}
\end{equation}
where $I_{\mathrm{h},n}$ is computed cheaply from the previous step's $\vv_n,\ii_n$.
$\mathbf{G}$ is factorised \emph{once}; each step is a single forward/back
substitution. Each element is a Norton equivalent (Norton's theorem),
\[
v = R_N\,(i + i_N), \qquad R_N = G^{-1},\;\; i_N \equiv \text{history current},
\]
and a distributed (Bergeron) line is two such Norton ends whose history currents
are the travelling waves from the \emph{other} end one propagation delay $\tau$
earlier,
\begin{equation}
I_{\mathrm{h},k}^{\text{line}}(t) = f\big(v_m,i_m;\,t-\tau\big),
\end{equation}
i.e. a delay rather than a derivative.

% ==========================================================================
\section{They are the same rule at different scopes}
% ==========================================================================

The two methods are not rival approximations: for a \textbf{linear network at a
fixed step} they produce the \emph{identical} sequence $y_n$, because both are
the trapezoidal discretisation of the same equations. The correspondence is
exact and visible in the coefficients. With $c_j = 2/\hstep$ from
\eqref{eq:yp},
\begin{equation}
G_L = \frac{\hstep}{2L} = \frac{1}{c_j\,L},
\qquad
G_C = \frac{2C}{\hstep} = c_j\,C .
\label{eq:equiv}
\end{equation}
\emph{The companion conductances of Method B are exactly the entries the
trapezoidal coefficient $c_j$ contributes to the Jacobian \eqref{eq:jac} of
Method A.} The history sources $I_{\mathrm{h},n}$ play the role of the constant term
$\alpha_n$ in \eqref{eq:yp}. The difference is therefore purely architectural:

\begin{itemize}
\item \textbf{Method A} substitutes \eqref{eq:trap} into the \emph{global
symbolic residual} and solves it \emph{implicitly} (Newton). Discretisation is a
property of the \textbf{solver}; models stay continuous.
\item \textbf{Method B} substitutes \eqref{eq:trap} into \emph{each element},
producing a constant real matrix and an explicit history, solved by one
\emph{linear} nodal pass. Discretisation is baked into the \textbf{models}.
\end{itemize}

% ==========================================================================
\section{Side-by-side}
% ==========================================================================

\begin{table}[h]
\centering
\renewcommand{\arraystretch}{1.25}
\begin{tabular}{>{\raggedright\arraybackslash}p{0.30\textwidth} >{\raggedright\arraybackslash}p{0.32\textwidth} >{\raggedright\arraybackslash}p{0.30\textwidth}}
\toprule
 & \textbf{A: Global DAE (Dynawo)} & \textbf{B: Companion (EMTP)} \\
\midrule
Element model & continuous, $v=L\dot i$ & pre-discretised, $i=G_L v+I_{\mathrm{h}}$ \\
Where \eqref{eq:trap} is applied & in the solver, globally & in each element, locally \\
Unknown each step & full state $y_{n+1}$ & node voltages $\vv_{n+1}$ \\
System solved each step & nonlinear $\Phi(y_{n+1})=0$ (Newton) & linear $\mathbf{G}\vv_{n+1}=\text{rhs}$ \\
Matrix & Jacobian $J$, rebuilt as needed & constant $\mathbf{G}$, factorised once \\
History kept by & solver ($y_n,\dot y_n$) & models ($I_{\mathrm{h}}$ via delay/sample) \\
Time step & variable or fixed & fixed (baked into $G_L,G_C$) \\
Nonlinearity / switching & inside Newton; refactor $J$ & rebuild/refactor $\mathbf{G}$; compensation \\
Cost per step & higher (nonlinear solve) & very low (one back-substitution) \\
Generality & any DAE, controls, phasor+EMT & EMT network, special-cased \\
\bottomrule
\end{tabular}
\caption{The opposing design choices. Both rest on the trapezoidal rule
\eqref{eq:trap}; \eqref{eq:equiv} shows their coefficients coincide.}
\end{table}

% ==========================================================================
\section{Consequences and trade-offs}
% ==========================================================================

\begin{description}
\item[Speed.] Method B is extremely fast: with a fixed step and fixed topology
the conductance $\mathbf{G}$ is factorised once and every step is a cheap
substitution. Method A pays for a nonlinear Newton solve over a larger system
each step.

\item[Robustness / generality.] Method A handles arbitrary nonlinearities,
stiff dynamics, controls, and even mixed phasor/EMT models uniformly, because
everything lives in one implicit residual and the integrator may adapt the step.
Method B special-cases nonlinear and switching elements (piecewise companions,
compensation/Newton sub-iterations) and a topology or step change forces a
re-factorisation.

\item[Numerical oscillations.] The trapezoidal rule is A-stable but not
L-stable: at a switching discontinuity it can ring (the well-known EMTP
``chatter''). Classical EMTP mitigates this with the \emph{Critical Damping
Adjustment} (half-step backward-Euler interpolation); a global-DAE solver
addresses the same effect through its event handling and step/order control.

\item[Step size.] Method B's accuracy and stability are tied to the single
fixed $\hstep$ chosen up front (typically $1$--$50\,\mu\mathrm{s}$).
Method A can in principle vary $\hstep$, though for waveform fidelity an EMT run
still uses a small step.
\end{description}

% ==========================================================================
\section{The implemented network primitives}
% ==========================================================================

Method A is abstract until one writes the residual each element actually
contributes to $F(y,\dot y,t)=0$. This section lists, for every
abc primitive of \texttt{DYNModelNetworkEMT}, (i) the residual as coded in its
\texttt{evalF}, and (ii) the classical companion counterpart, so the two columns
of the side-by-side table become concrete per device.

\paragraph{Conventions.} Receptor convention; instantaneous \emph{peak}-pu abc
quantities, $v_k=\sqrt2\,V_{\mathrm{rms}}\cos(\theta-2\pi k/3)$, $k\in\{a,b,c\}$.
From the IIDM data: $R=r/Z_{\mathrm b}$, $L=x/(Z_{\mathrm b}\,\omega_N)$,
$C=b\,Z_{\mathrm b}/\omega_N$, with $\omega_N=2\pi f_N$. Below, $v_k,i_k$ are the
node-voltage / branch-current entries of $y$ owned by the component, and
$\dot{(\cdot)}\equiv\dd(\cdot)/\dd t$ is the derivative the solver discretises
through \eqref{eq:yp}.

\subsection{Node (bus): \texttt{ModelBusEMT}}
A bus is a small shunt capacitance $C$ to ground (physical/stray $C$, plus any
shunt-compensator $C$) into which every incident branch injects its current
(KCL), with shunt conductance $G$ (loads) and an optional fault conductance
$G_f$. Per phase $k$ the residual is
\begin{equation}
f_k \;=\; \sum_{\text{branches}}(\pm i_k)\;-\;(G_k+G_{f,k})\,v_k\;-\;C\,\dot v_k
\;=\;0,
\qquad\text{i.e.}\quad C\,\dot v_k=\textstyle\sum i_k-(G_k+G_{f,k})v_k .
\label{eq:bus}
\end{equation}
It is differential whenever $C>0$; the network enforces $C\ge 10^{-4}$ (a snubber)
so every node is a stiff \textbf{index-1} differential state --- the DAE
regularisation that takes the place of EMTP's ever-present $G_C=2C/\hstep$. A
dead (switched-off) island pins $f_k=v_k$ ($v=0$).
\emph{Companion (EMTP):} the capacitor becomes $G_C=2C/\hstep \parallel
I_{\mathrm h}^C$ and the node row reads $(G_C+\sum G)\,v_k=I_{\mathrm h}^C+\sum
i^{\text{hist}}_{\text{branch}}$; a $C$-less node is simply a conductance row.

\subsection{Series R--L branch (line, cable, reactor): \texttt{ModelLineEMT}}
The branch current $i_k$ is a state and the constitutive law is the residual,
\begin{equation}
f_k \;=\; v_{\mathrm{from},k}-v_{\mathrm{to},k}-R\,i_k-L\,\dot i_k \;=\;0,
\label{eq:line}
\end{equation}
with $-i_k$ injected into the \emph{from} node and $+i_k$ into the \emph{to} node
(the $\sum i_k$ of \eqref{eq:bus}); open $\Rightarrow f_k=i_k$. This is
\eqref{eq:Aexample} embedded in the network. \emph{Companion (EMTP):} the same
branch is the conductance $G_{RL}=1/(R+2L/\hstep)$ in parallel with a history
source, the current being back-substituted after the nodal solve rather than
carried as a state --- the single sharpest contrast between the two methods.

\subsection{Switch / breaker: \texttt{ModelSwitchEMT}}
An ideal switch is a pure constraint on $y$:
\begin{equation}
\text{closed: } f_k=v_{1,k}-v_{2,k}=0,
\qquad
\text{open: } f_k=i_k=0,
\label{eq:switch}
\end{equation}
the current unknown $i_k$ (injected $\mp1$ into the two nodes) absorbing the
zero-impedance equality. \emph{Companion (EMTP):} an ideal zero-impedance switch
makes the nodal matrix singular, so EMTP uses a small $R_{\mathrm{on}}$ / large
$R_{\mathrm{off}}$ or node merging; the augmented DAE set represents the ideal
constraint directly.

\subsection{Two-winding transformer: \texttt{ModelTwoWindingsTransformerEMT}}
With turns ratio $r_T=N_1/N_2$ and phase shift $\alpha$ (an abc rotation
$\mathcal R(\alpha)$), the primary current $i_k$ is the state:
\begin{equation}
f_k=v_{\mathrm{pri},k}-r_T\,[\mathcal R(\alpha)\,v_{\mathrm{sec}}]_k-R\,i_k-L\,\dot i_k=0,
\quad
[\mathcal R(\alpha)v]_k=\cos\alpha\,v_k+\tfrac{\sin\alpha}{\sqrt3}\big(v_{k+2}-v_{k+1}\big),
\label{eq:tfo}
\end{equation}
with $-i_k$ injected into the primary node and the power-conserving
$r_T\,[\mathcal R(-\alpha)\,i]_k$ into the secondary. For $r_T=1,\alpha=0$ it
reduces to the R--L branch \eqref{eq:line}. \emph{Companion (EMTP):} an
ideal-transformer (turns-ratio) companion in series with the R--L companion of
\eqref{eq:line}; the rotation $\mathcal R(\alpha)$ is identical.

\subsection{Default load and PQ generator: \texttt{ModelLoadEMT}, \texttt{ModelGeneratorEMT}}
A constant-impedance load owns no state and registers a per-phase conductance
$G=1/R$ into its node (the $G$ of \eqref{eq:bus}) --- algebraic and identical in
both methods. The default generator is a prescribed balanced current source
\begin{equation}
i_k(t)=\sqrt2\,I\,\cos\!\big(\omega_N t+\varphi-\tfrac{2\pi k}{3}\big),
\label{eq:gen}
\end{equation}
injected into the node --- again the same object (a current source) in both
methods. Voltage-dependent / restorative variants are deliberately \emph{not}
network primitives: their defining quantity is the RMS $U$, a state in the
quasi-static phasor world but a \emph{measurement} in EMT, so they belong to a
Modelica model on the bus terminal, not here.

\subsection{Synchronous machine: \texttt{Electrical/EMT/GeneratorSynchronous}}
The detailed machine is a Modelica model bound to its bus through the abc ACPIN.
It is a Park ($dq0$) model that, unlike the phasor, \textbf{keeps the stator
electromagnetic transients} $\dot\lambda/\omega_N$. The abc terminal is mapped to
$dq0$ by Park's transform,
\begin{equation}
u_d=C\!\sum_k\sin\!\big(\theta-\tfrac{2\pi k}{3}\big)v_k,\quad
u_q=C\!\sum_k\cos\!\big(\theta-\tfrac{2\pi k}{3}\big)v_k,\quad
u_0=\tfrac{1}{\sqrt3}\textstyle\sum_k v_k,
\label{eq:park}
\end{equation}
and the stator current is injected back through the inverse map
$i_k=K_i C_{dq}\big[\sin(\theta-\tfrac{2\pi k}{3})i_d+\cos(\theta-\tfrac{2\pi
k}{3})i_q+i_0/\sqrt2\big]$. The stator voltage equations \emph{retain} the
transient term (the EMT difference --- the phasor sets $\dot\lambda=0$):
\begin{align}
u_d&=(R_a+R_{Tfo})\,i_d+\dot\lambda_d/\omega_N-\omega\,\lambda_q,\nonumber\\
u_q&=(R_a+R_{Tfo})\,i_q+\dot\lambda_q/\omega_N+\omega\,\lambda_d,\qquad
u_0=R_a\,i_0+\dot\lambda_0/\omega_N,
\label{eq:stator}
\end{align}
the rotor windings ($f,D,Q_1,Q_2$) obey $u_f=R_f i_f+\dot\lambda_f/\omega_N$,
$0=R_D i_D+\dot\lambda_D/\omega_N,\dots$, the fluxes link the currents through the
saturable mutuals $M_{d,\mathrm{sat}},M_{q,\mathrm{sat}}$ (Canay $M_{rc}$), e.g.
$\lambda_d=(M_{d,\mathrm{sat}}+L_d+X_{Tfo})i_d+M_{d,\mathrm{sat}}i_f+M_{d,\mathrm{sat}}i_D$,
$\lambda_0=L_0 i_0$, and the swing closes the loop,
\begin{equation}
\dot\theta=\omega\,\omega_N,\qquad
2H\,\dot\omega=c_m\tfrac{P_{\text{NomTurb}}}{S_{\text{Nom}}}-c_e-D_{Pu}(\omega-\omega_{\text{ref}}),\qquad
c_e=\lambda_q i_d-\lambda_d i_q .
\label{eq:swing}
\end{equation}
\emph{Companion (EMTP):} a universal-machine companion couples to the network by
a Norton/Th\'evenin source at the stator terminal (an internal EMF behind a
subtransient impedance, predicted and updated each fixed step). Dynawo instead
solves the machine fluxes \eqref{eq:stator}--\eqref{eq:swing} and the network abc
voltages \eqref{eq:bus} \emph{simultaneously} as one DAE --- the ACPIN equates
$\texttt{terminal.v}\leftrightarrow$ node $v$ and
$\texttt{terminal.i}\leftrightarrow$ the KCL injection --- so there is no
predictor/companion lag at the machine--network interface.

\subsection{Summary}
\begin{table}[h]
\centering
\renewcommand{\arraystretch}{1.25}
\begin{tabular}{>{\raggedright\arraybackslash}p{0.20\textwidth}
  >{\raggedright\arraybackslash}p{0.18\textwidth}
  >{\raggedright\arraybackslash}p{0.52\textwidth}}
\toprule
Primitive & State(s) & Implemented residual / contribution \\
\midrule
Bus & $v_k$ (if $C{>}0$) & $C\dot v_k=\sum i_k-(G{+}G_f)v_k$ \eqref{eq:bus} \\
R--L branch & $i_k$ & $v_{\mathrm{from}}-v_{\mathrm{to}}-Ri_k-L\dot i_k=0$ \eqref{eq:line} \\
Switch & $i_k$ & closed $v_1{-}v_2=0$, open $i_k=0$ \eqref{eq:switch} \\
2W transformer & $i_k$ & $v_{\mathrm{pri}}-r_T\mathcal R(\alpha)v_{\mathrm{sec}}-Ri_k-L\dot i_k=0$ \eqref{eq:tfo} \\
Load & --- & shunt $G=1/R$ into the node \\
PQ generator & --- & prescribed current source \eqref{eq:gen} \\
Machine & $\lambda_{d,q,0,f,D,Q_1,Q_2},\theta,\omega$ & Park $dq0$ with $\dot\lambda/\omega_N$ \eqref{eq:park}--\eqref{eq:swing} \\
\bottomrule
\end{tabular}
\caption{The abc EMT primitives and the residuals they contribute to the global
DAE $F(y,\dot y,t)=0$. Every $\dot{(\cdot)}$ is discretised by the solver through
\eqref{eq:yp}; the classical method would replace each storage element by its
companion model instead.}
\end{table}

% ==========================================================================
\section{The Modelica EMT component library}
% ==========================================================================

Dynawo carries EMT primitives at \emph{two} layers, both in the continuous
(Method A) form. Section~7 covered the C++ \texttt{DYNModelNetworkEMT}, which
builds an abc network directly from the IIDM. Alongside it lives a standalone
\textbf{Modelica} library, \texttt{Dynawo.Electrical.EMT}, of preassembled abc
components that are \emph{composed through \texttt{.dyd} connections} (or attached
to the C++ network's ACPIN). The two share conventions and waveforms; this
section records the Modelica components whose formulation is not already a C++
primitive.

\paragraph{Two-terminal base.} Every passive element extends
\texttt{BaseEmtTwoTerminal}, which declares the branch voltage and the
flow-connector KCL
\begin{equation}
v_k = p.v_k - n.v_k,\qquad p.i_k + n.i_k = 0,\qquad k\in\{a,b,c\},
\label{eq:base2t}
\end{equation}
($p.i$ the through-current; the \texttt{flow} connector sums currents at shared
nodes). The constitutive laws are then literal continuous equations:
\begin{equation}
\text{R: } v_k=R\,i_k,\quad
\text{L: } v_k=L\,\dot i_k,\quad
\text{C: } i_k=C\,\dot v_k,\quad
\text{SeriesRL: } v_k=R\,i_k+L\,\dot i_k .
\label{eq:rlc}
\end{equation}

\paragraph{Coupled R--L (mutual phases): \texttt{CoupledRL}.} The one passive
element \emph{not} diagonal --- the abc series impedance with inter-phase mutuals
written as $3\times3$ matrices,
\begin{equation}
\mathbf v = \mathbf R\,\mathbf i + \mathbf L\,\dot{\mathbf i},
\qquad \mathbf R,\mathbf L\in\mathbb R^{3\times3},
\label{eq:coupledrl}
\end{equation}
the full mutual-coupling generalisation of the diagonal $\texttt{SeriesRL}$ /
C++ line \eqref{eq:line}.

\paragraph{Winding-connection transformers.}
\texttt{TransformerYgYg} (grounded-wye / grounded-wye, $1{:}1$ phase) is the
per-phase
\begin{equation}
v_{1,k}=r_T\,v_{2,k}+R\,i_{1,k}+L\,\dot i_{1,k},
\qquad i_{2,k}=-r_T\,i_{1,k},
\label{eq:ygyg}
\end{equation}
i.e.\ the C++ branch \eqref{eq:tfo} at $\alpha=0$. \texttt{TransformerYgD}
(grounded-wye / delta) realises the $30^\circ$ shift \emph{through the winding
connection} rather than a rotation matrix: secondary winding $k$ sees the delta
difference $v_{2,k}-v_{2,k+1}$, and the line current is the difference of two
winding currents,
\begin{equation}
v_{1,k}=r_T\big(v_{2,k}-v_{2,k+1}\big)+R\,i_{1,k}+L\,\dot i_{1,k},
\qquad
i_{2,k}=-r_T\big(i_{1,k}-i_{1,k-1}\big),
\label{eq:ygd}
\end{equation}
(indices mod 3). The $\sqrt3$ magnitude and $30^\circ$ phase fall out of the
delta topology --- the physical EMT way to get the shift the C++ branch
approximates with $\mathcal R(\alpha)$.

\paragraph{Ideal sources.} \texttt{VoltageSource} imposes an ideal abc EMF
$v_k=\hat U\cos(\Omega t+\varphi+\psi_k)$, $\psi_k\in\{0,-\tfrac{2\pi}{3},
+\tfrac{2\pi}{3}\}$. \texttt{InfiniteBus} is an EMF behind a small series $R_s$,
which makes it connectable to any (inductive or capacitive) node:
\begin{equation}
i_k=\frac{v_k-\sqrt2\,U\cos(\Omega t+\varphi+\psi_k)}{R_s},
\qquad \Omega=2\pi f_N .
\label{eq:infbus}
\end{equation}

\paragraph{Faults.} \texttt{Fault} is the (multi-)phase-to-ground short: while
active, each faulted phase obeys $v_k=R_f\,i_k$ (small $R_f$); every healthy phase
draws the negligible $i_k=v_k/R_\infty$. \texttt{LineToLineFault} shorts two
phases through $R_f$,
\begin{equation}
v_{p_1}-v_{p_2}=R_f\,i_F,\qquad i_{p_1}=i_F,\quad i_{p_2}=-i_F,\quad
i_{p_3}=v_{p_3}/R_\infty,
\label{eq:llfault}
\end{equation}
the abc analogue of the C++ network's $G_f$ shunt-to-ground term in
\eqref{eq:bus} (here written as a Modelica terminal model). \texttt{LineFault}
splices such a fault into a line.

\paragraph{Machines.} The three Modelica machines all use the Park $dq0$ +
$\dot\lambda/\omega_N$ stator form of \eqref{eq:park}--\eqref{eq:swing};
they differ in detail, not in method:
\texttt{GeneratorSynchronous} (Section~7.6) carries an integrated step-up
transformer ($L_d{+}X_{Tfo}$, $R_a{+}R_{Tfo}$) and an \texttt{omegaRefPu} input;
\texttt{SynchronousMachineFull} is its standalone twin in the \emph{absolute}
rotor frame ($\dot\theta=\omega\,\omega_N$, no \texttt{omegaRef}, swing damping
$D_{Pu}(\omega-1)$), with saturation evaluated on the \emph{instantaneous} air-gap
flux (singularity-guarded) --- the physically correct behaviour under unbalance /
faults; \texttt{SynchronousMachine} is the reduced two-axis (subtransient) model.
All keep the EMT zero-sequence winding $\lambda_0=L_0 i_0$.

% ==========================================================================
\section{A concrete Method-B reference: classical EMTP travelling-wave lines}\label{sec:emtp-lines}
% ==========================================================================

The companion/nodal Method~B of Section~3 has a concrete counterpart in the
classical EMTP family of tools, sketched here purely as the \emph{opposite} of the
Dynawo DAE primitives of Sections~7--8. A first observation: most \emph{lumped}
EMTP elements are themselves written in the continuous $\dot{}$ form of
Method~A --- an R-L branch is literally $v=R\,i+L\,\dot i$, and transformers and
synchronous machines are continuous too. The genuinely Method-B part --- the part
that has \emph{no} Dynawo equivalent --- is the \textbf{travelling-wave
transmission lines}.

\paragraph{Norton end.} Every line terminal is a Norton companion:
\begin{equation}
v = R_N\,(i + i_N),\qquad \textstyle\sum \texttt{pin.i}=0,
\label{eq:emtp-norton}
\end{equation}
a fixed conductance $G_N=1/R_N$ in parallel with a \emph{history current source}
$i_N$ --- exactly the Method-B element \eqref{eq:compL}--\eqref{eq:nodal}, now a
concrete model.

\paragraph{Constant-parameter (Bergeron) line.} For a line of characteristic
impedance $Z_c$ and one-way travel time $\tau$, $R_N=Z_c$ and the two ends couple
\emph{only} through the wave that left the far end one delay ago. In
the history term this is literally a \texttt{delay}:
\begin{equation}
i_{h,k}(t) = \texttt{delay}\big(i_{r,m},\,\tau\big),\qquad
i_{h,m}(t) = \texttt{delay}\big(i_{r,k},\,\tau\big),
\label{eq:emtp-bergeron}
\end{equation}
($k,m$ the two ends). The history is a \textbf{look-up of past values}, not a
derivative of present unknowns; and because $\tau\ge\hstep$ the two end nodes are
\emph{decoupled within a step} --- the structural reason EMTP's $\mathbf G$ is
sparse/block-diagonal and the per-step solve is cheap and parallel.

\paragraph{Wideband (frequency-dependent) line.} The same Norton ends, but the
history now carries the line's frequency dependence as a recursive convolution
(a pole--residue fit of the modal propagation/admittance), realised by auxiliary
history states:
\begin{equation}
\dot x_{k,p,i}=P_{p}\,x_{k,p,i}+R_{p,i}\,i_{h,m},
\qquad
i_{N,k}=\sum_{p,i} x_{k,p,i},
\label{eq:emtp-wb}
\end{equation}
i.e.\ $i_N(t)=\int a(t-s)\,i(s)\,\dd s$ approximated by first-order states.

\paragraph{Why this is the opposite of the DAE.} A Dynawo line is the single
implicit equation $v_{\mathrm{from},k}-v_{\mathrm{to},k}-R\,i_k-L\,\dot i_k=0$
\eqref{eq:line} spanning \emph{both} ends and the current state, solved
\emph{simultaneously} with the whole network by Newton each step. An EMTP
travelling-wave line is \emph{two} Norton sources \eqref{eq:emtp-norton} whose
only link is a \texttt{delay} \eqref{eq:emtp-bergeron} (or a convolution
\eqref{eq:emtp-wb}) of \emph{past} values. The contrast is total and is the whole
point of the note:

\begin{center}\renewcommand{\arraystretch}{1.25}
\begin{tabular}{>{\raggedright\arraybackslash}p{0.46\textwidth}>{\raggedright\arraybackslash}p{0.46\textwidth}}
\toprule
\textbf{Dynawo line (DAE, \S7)} & \textbf{EMTP travelling-wave line (companion)} \\
\midrule
$v_{\mathrm{from}}-v_{\mathrm{to}}-R i-L\,\dot i=0$ & two Norton ends $v=R_N(i+i_N)$ \\
ends coupled \emph{implicitly}, in one equation & ends coupled only by $\texttt{delay}(\cdot,\tau)$ / convolution \\
$\dot i$ --- derivative of present unknown & $i_N$ --- history of past values \\
solved with the whole network (Newton) & cheap decoupled nodal solve per step \\
\bottomrule
\end{tabular}
\end{center}

Feeding such a line to a DAE solver is the category error noted in the next
section: its \texttt{delay()} history is an explicit past-value look-up, foreign
to a continuous Newton step --- so a port to Dynawo replaces the travelling-wave
companion with a lumped, continuous $\dot{}$ segment and lets the solver own the
time stepping.

% ==========================================================================
\section{Mapping to this repository}
% ==========================================================================

\begin{itemize}
\item \textbf{Method A} is what this repository implements: Dynawo's
\texttt{SolverTRAP} (\eqref{eq:yp}) together with the instantaneous-abc
primitives of the C++ \texttt{DYNModelNetworkEMT} (Section~7) and the Modelica
\texttt{Dynawo.Electrical.EMT} library (Section~8), each written in continuous
$\dot{}$ form (e.g.\ \texttt{SeriesRL}: $v=R\,i+L\,\dot i$). The validated demo
\texttt{ThreePhaseRLC} matches the analytic steady state to $\le 0.1\%$. Every
EMT model named in this note is a Dynawo model.

\item \textbf{Method B} (the classical companion/nodal form) is \emph{not} a
Dynawo model; it is dissected in Section~\ref{sec:emtp-lines}: the EMTP
travelling-wave lines (Norton ends coupled by a \texttt{delay} or a convolution
history term) are the genuine companion form, while the lumped $RLC$ /
transformer / machine models are already continuous (Method~A).

\item \textbf{Practical note.} Feeding a companion model (Method B) to a DAE
solver (Method A) is a category error: the element is already discretised, and
its delayed/sampled history fights the continuous Newton step. A model written
for Method A keeps the continuous $\dot{}$ form and lets the solver, not the
model, own the time stepping.
\end{itemize}

\end{document}
